\chapter{Integrated Discussion and Findings}

This chapter brings together what was learned across PBL I, II and III. Numbers are compared only between models
evaluated on the same task, the same test set and the same metric. Results from different datasets are compared
qualitatively.

\section{Summary of Outcomes per Unit}
\begin{table}[H]
\centering
\caption{Best result of each unit on its own held-out test data}
\label{tab:summary}
\renewcommand{\arraystretch}{1.25}
\begin{tabularx}{\linewidth}{@{}c >{\raggedright\arraybackslash}p{3.3cm} >{\raggedright\arraybackslash}p{3.4cm} >{\raggedright\arraybackslash}X@{}}
\toprule
\textbf{Unit} & \textbf{Task / dataset} & \textbf{Best method} & \textbf{Key result (vs baseline)} \\
\midrule
1 & EDA: Play Store, Video Games & Robust cleaning + rank statistics & Quantified market drivers (e.g.\ reviews–installs $\rho=0.97$; top 1\% of games = 21\% of sales) \\
2 & Regression: CCPP & Tuned XGBoost & RMSE 3.17 MW, $R^2$ 0.966 (linear: 4.59 MW; $-$31\% error) \\
2 & Classification: Bank Marketing & Weighted XGBoost + threshold & F1 0.531, ROC-AUC 0.812, 3.3$\times$ lift (plain LR: F1 0.320) \\
3 & Clustering: UCI HAR & Gaussian Mixture & Accuracy 0.612, ARI 0.486 (K-Means: 0.570) \\
4 & Sequential: UCI HAR & Hybrid HMM + Viterbi & Accuracy 0.987, macro F1 0.988 (same classifier without HMM: 0.964) \\
5 & Ensembles: Diabetes & \dbBestName{} + tuned threshold & F1 \dbBestFone, ROC-AUC \dbBestAuc{} (LR: F1 \dbLrFone, ROC-AUC \dbLrAuc) \\
\bottomrule
\end{tabularx}
\end{table}

\section{Critical Findings Across the Project}
\begin{enumerate}
\item \textbf{Combining many learners beats a single model on tabular data, consistently.} Gradient-boosted tree
ensembles were the best model on every supervised tabular task: CCPP regression (tuned XGBoost), bank-marketing
classification (weighted XGBoost) and diabetes prediction (Chapter 4, where every one of the ten ensembles beat both
baselines on F1 and ROC-AUC). The mechanism is the same in each case. Boosting captures non-linear, interacting effects
that linear models miss, and averaging many trees removes the instability of a single tree.

\item \textbf{The temporal structure adds the same kind of value as an ensemble.} In Unit 4 a strong per-window
classifier made 3.6\% errors. Adding a Markov chain and decoding with Viterbi removed 63\% of them, because each
prediction is combined with the evidence from neighbouring windows. Like an ensemble, the HMM pools evidence to correct
isolated mistakes; here the pooling runs over time rather than over models.

\item \textbf{The right metric matters more than the model when classes are imbalanced.} Three datasets were
imbalanced: bank 11.3\% positive, diabetes 8.8\%, and the activity classes once sequences are considered. On each, a
trivial or plain model reached high accuracy while missing most positives (bank: 89.8\% accuracy with only 21\% recall).
Evaluating with F1, recall, ROC-AUC and PR-AUC, and choosing the decision threshold on out-of-fold training predictions,
consistently gave more useful models: bank F1 0.487 $\rightarrow$ 0.531, and a similar threshold gain for diabetes
(Chapter 4).

\item \textbf{Leakage control is essential for honest results.} Each review had its own leakage risk: the
\texttt{duration} feature in the bank data, duplicate patient records in the diabetes data, and person overlap in HAR.
Each was removed or controlled (feature dropped, duplicates removed before splitting, split by subject). Preprocessing
and model selection always ran inside training folds.

\item \textbf{Unsupervised structure is not the same as semantic labels.} K-Means clusters of HAR windows were
geometrically \emph{tighter} than the true activities (silhouette 0.41 vs 0.26) yet matched them only 57\% of the time.
Unsupervised methods reveal how data is organised, which is valuable for exploration, but they do not replace labels
when a specific target matters.

\item \textbf{Interpretation confirms domain knowledge.} Permutation importance and SHAP identified ambient temperature
as the driver of power output, past campaign success and the macro-economy as drivers of subscription, and HbA1c,
glucose, age and BMI as drivers of diabetes risk. In each case the model's explanations matched the physics, the
marketing literature or clinical guidelines, which increases trust in the models.
\end{enumerate}

\section{Comparability of Results}
Only models on the same dataset, split and metric are compared numerically: the rows of each table in Chapters 2–4.
The six summary rows of Table~\ref{tab:summary} come from different datasets with different difficulty and class
balance. For example, the diabetes F1 (\dbBestFone) and the bank F1 (0.531) should not be read as ``the diabetes model is
better''. The diabetes data contains much stronger predictors (HbA1c and glucose), while the bank data is limited to
information available before a call.

\section{Study Limitations}
\begin{itemize}
\item \textbf{Different datasets per unit.} The allotted topics used different datasets, so the five units cannot be
evaluated on one common benchmark.
\item \textbf{Public, curated data.} All datasets are public snapshots: laboratory HAR recordings, a 2008–2010 bank
campaign, and a curated diabetes dataset whose laboratory values are recorded at a few discrete levels. Performance on
raw clinical or field data is likely to be lower.
\item \textbf{Tuning budget.} Only the leading family in each task was tuned extensively; the other models used
sensible defaults.
\item \textbf{No external validation.} Each model was validated on a held-out split of the same source, not on an
independent cohort or site.
\end{itemize}
